\documentstyle[% 


righttag% 


,11pt% 


,amssymb% 


,russian%               remove in Eng. 


%,izvru%                 Set izven% in Eng.


,izvru-ks%             for brief communications


]{amsart} 


 


% 


\newcommand{\tts}[1]{}  


\renewcommand{\baselinestretch}{0.96}


 


%\hoffset=-15mm%                  For Eng. 


\setcounter{page}{81} 


\god{2002} 


\nomer{7~(482)} %                        Remove in Eng. 


%\nomer{Vol.\,46, No.\,7}%               Set in Eng. 


%\str{\pageref{begin}--\pageref{end}}%  Set in Eng. 


\udk{519.854.2} 


 


\author{\it А.Б.\,Зинченко} 


% NB:   ^^ remove in Eng. 


\title{Структура множества Парето некоторых векторных задач  


составления расписаний} 


 


\date{} 


%\pagestyle{myheadings}%         Set in Eng. 


\pagestyle{plain} %              Remove in Eng. 


\begin{document} 


%\thispagestyle{plain}%          Set in Eng. 


\maketitle 


\label{begin} 


 


 


{\bf 1.} {\it Постановка задачи.} Имеется $n\ge2$ независимых работ, для 


выполнения 


которых требуется $m\ge1$ типов ресурсов. Интенсивности выделения 


ресурсов постоянны во времени и ограничены величинами $\beta_j>0$, 


$j\in M=\{1,\dotsc,m\}$. Интенсивности потребления ресурсов работами 


также постоянны. 


Известны длительности $a_i>0$, $i\in N=\{1,\dotsc,n\}$, выполнения работ и 


ресурсные 


потребности $0\le b_{i j}\le\beta_j$, $i\in N$, $j\in M$. Требуется 


составить расписание 


без прерываний, по возможности минимизирующее время завершения всех 


%% tts comma has been placed  ^^ here previously; the sense is violated


работ, а также необходимое количество каждого ресурса. Рассматриваемая 


задача является $NP$-трудной, т.\,к. известная $NP$-полная проблема 


упаковки в 


контейнеры ([1], с.\,20) является частным случаем более простой 


(одноресурсной и однокритериальной) задачи, описанной в [2]. 


 


Пусть $t_i$ --- момент завершения работы $i$. Введем инъективное 


отображение $I$, сопоставляющее каждому расписанию $t\in T_1=\{t\in 


R^n:t\ge a\}$ 


систему $\{I_1(t),\dotsc,I_n(t)\}$ открытых интервалов 


$I_i(t)=(t_i-a_i,t_i)$, $i\in N$, вещественной 


оси. Определим множество подмножеств индексов пересекающихся 


интервалов $\xi(I(t))=(\mathop{\cup}\limits_{i\in N}\{i\})\bigcup 


\{e\in2^N:|e|\ge2,\ \mathop{\cap}\limits_{i\in e}I_i(t)\ne\emptyset\}$. 


 


Выразив длину расписания $t$ и общие ресурсные потребности работ 


при расписании $t$ в виде функций 


$$ 


f_0(t)=\max_{i\in N}t_i;\quad 


f_j(t)=\max_{e\in\xi(I(t))}v_j(e),\ \ j\in M, 


$$ 


где $v_j(e)=\sum\limits_{i\in e}b_{i j}$, $j\in M$; 


$e\in2^n\setminus\emptyset$, получаем невыпуклую векторную задачу 


\begin{equation} 


f(t)=(f_0(t),\dotsc,f_m(t))\to\min,\quad 


t\in T_2=\{t\in T_1:f_j(t)\le\beta_j,\ j\in M\}, 


\end{equation} 


которая может иметь бесконечное множество эффективных (оптимальных по 


Парето) решений $\bold P_T$. Множество эффективных оценок $\bold P_f$ 


задачи (1), как 


показано ниже, всегда конечно. 


%\smallskip 


 


{\bf 2.} {\it Свед\'ение к траекторной задаче.} Пусть $U=\{(i,j):i,j\in 


N,\ i<j\}$, a $G$ --- 


биекция из $2^U$ на множество простых $n$-вершинных неориентированных 


графов $(N,E)$, $E\in2^U$. Рассмотрим отображение $\Psi$, которое каждой 


системе 


интервалов $I(t)$ соотносит множество пар индексов пересекающихся 


интервалов, т.\,е. $\Psi(I(t))=\{(i,j)\in U:I_i(t)\cap 


I_j(t)\ne\emptyset\}$. Тогда $G\circ\Psi\circ I$ сюръективно 


отображает множество расписаний $T_1$ на множество $n$-вершинных 


графов, представимых интервалами, длины которых равны длительностям 


выполнения работ. Обозначим $\Omega_0=\{E\in2^U:G(E)$ 


--- граф интервалов$\}$, $\Omega_1=\Psi(I(T_1))\subseteq\Omega_0$. 


Представив $T(E)=(\Psi\circ I)^{-1}(E)$ в виде $T(E)=X(E)\cap Y(E)$, где 


$$ 


X(E)=\{t\in T_1:I_i(t)\cap I_j(t)=\emptyset, 


\ (i,j)\in U\setminus E\},\quad 


Y(E)=\{t\in R^n:I_i(t)\cap I_j(t)\ne\emptyset, 


\ \ (i,j)\in E\}, 


$$ 


получим разбиение множества расписаний 


$T_1=\mathop{\cup}\limits_{E\in\Omega_1}T(E)$. 


 


Обозначим через $K(E)$ ($H(E)$) семейство кликовых (независимых) 


множеств графа $G(E)$. Будем рассматривать вершинно-взвешенные графы, 


т.\,е. графы, $i$-й вершине которых приписаны числа 


$(a_i,b_{i1},\dotsc,b_{i m})$. На 


множестве $2^U$ зададим функции 


$$ 


F_0(E)=\max_{e\in H(E)}v_0(e),\quad 


F_j(E)=\max_{e\in K(E)}v_j(e),\ \ j\in M, 


$$ 


где $v_0(e)=\sum\limits_{i\in e}a_i$, $e\in2^N\setminus\emptyset$. 


Функция $F_0$ ($F_j$, $j\ne0$) антитонна (изотонна) 


на $(2^U,\subseteq)$, т.\,е. $E^1\subseteq E^2\Rightarrow F_0(E^1)\ge 


F_0(E^2)$; $F_j(E^1)\le F_j(E^2)$, $j\in M$. 


Для любого $E\in\Omega_0$ значения $F_0(E)$ и $F_j(E)$, $j\in M$, 


полиномиально вычисляемы [3], [4]. 


 


Очевидно, для любых $E\in\Omega_1$, $t\in T(E)$ справедливы равенства 


$f_j(t)=F_j(E)$, $j\in M$, из которых следует разложение множества 


допустимых 


расписаний $T_2=\mathop{\cup}\limits_{E\in\Omega_1\cap\Omega_2} T(E)$, 


где $\Omega_2=\{E\in2^U:F_j(E)\le\beta_j,\ j\in M\}$. 


 


Tак как функции $f_j$, $j\in M$, постоянны на множестве $T(E)$, то 


парето-оптимальными могут быть только точки минимума функции $f_0$ на 


$T(E)$. 


 


Пусть $E\in\Omega_1$ и $t\in T(E)$, тогда для любого $h(E)\in H(E)$ 


интервалы 


$I_j(t)$, $j\in h(E)$, попарно не пересекаются. Следовательно, 


$f_0(t)\ge F_0(E)$, т.\,е. 


верхняя оценка функции $f_0$ на $T(E)$ определяется независимым множеством 


максимального веса графа $G(E)$. 


 


Для описания структуры множества Парето $\bold P_T$ задачи (1) введем 


вспомогательную дискретную задачу: 


\begin{equation} 


F(E)=(F_0(E),\dotsc,F_m(E))\to\min,\quad 


E\in\Omega_3=\{E\in\Omega_1\cap\Omega_2:T^*(E)\ne\emptyset\}, 


\end{equation} 


где $T^*(E)=\{t\in T(E):f_0(t)=F_0(E)\}$. 


 


\begin{thm} 


Справедливо разложение 


\begin{equation} 


\bold P_T=\bigcup_{E\in P_\Omega}T^*(E), 


\end{equation} 


где $\bold P_\Omega$ --- множество Парето задачи $(2)$. 


\end{thm} 


 


\begin{cor} 


Задача (1) сводима к задаче (2). 


\end{cor} 


 


\begin{cor} 


Если $Q_1,\dotsc,Q_q$ --- классы эквивалентности множества 


$\bold P_\Omega$ относительно критерия $F$, то 


$\mathop{\cup}\limits_{E\in Q_1}T^*(E),\dotsc, 


\mathop{\cup}\limits_{E\in Q_q}T^*(E)$ --- классы 


эквивалентности 


множества $\bold P_T$ относительно критерия $f$. 


\end{cor} 


 


\begin{cor} 


Справедлива достижимая оценка $|\bold P_f|\le2^{n(n-1)/2}$, где $\bold 


P_f=\{f(t):t\in\bold P_T\}$. 


\end{cor} 


 


{\bf 3.} {\it Нахождение эффективных решений.} Задача (2) относится к 


классу 


траекторных [5]. Траекториями являются множества ребер $n$-вершинных, 


$(m+1)$-взвешенных интервальных графов, обладающих определенными 


свойствами. 


Множество траекторий есть нижняя подполурешетка булевой 


решетки $(2^U,\subseteq)$. Частные критерии монотонны и имеют тип 


MINMAXSUM. 


 


Эффективные траектории можно выделить комбинаторными алгоритмами, 


краткий обзор которых дан в ([6], с.\,186). Эти алгоритмы оcнованы 


на погружении допустимой области в множество простой структуры 


и требуют описания процедур, проверяющих допустимость точки ($g$-оракул) 


и вычисляющих значение критерия ($F$-оракул). Поскольку задача (2) имеет 


критериальные ограничения, то $g$-оракул одновременно является и 


$F$-оракулом. 


Назовем его $F,g$-оракулом. 


%\pagebreak


 


\begin{center} 


{\bf Алгоритм 1} ($F,g$-оракул задачи (2)) 


\end{center} 


 


1. Дано $E\in2^U$. Если $T(E)=\emptyset$, то $E\notin\Omega_3$; 


перейти на п.\,5. 


 


2. Вычислить $F_0(E)$. Если $T^*(E)=\emptyset$, то $E\notin\Omega_3$; 


перейти на п.\,5. 


 


3. Проверить условия $F_j(E)\le\beta_j$, $j\in M$. Если хоть одно из 


них не 


выполняется, то $E\notin\Omega_3$; перейти на п.\,5. 


 


4. Получили $E\in\Omega_3$. 


 


5. Конец. 


%\smallskip 


 


Временн\'ая сложность алгоритма определяется трудоемкостью 


проверки условия $T^*(E)=\emptyset$. Множество $T^*(E)$ задается системой, 


включающей 


дизъюнктивные неравенства, определяющие $X(E)$, строгие линейные 


неравенства, 


определяющие $Y(E)$, и уравнение $f_0(t)=F_0(E)$. Эффективные алгоритмы 


решения систем такого типа не известны [7]. 


 


Задача несколько упрощается, если находить не все множество $\bold P_T$, а 


только полное множество альтернатив $\bold P^0_T$. 


 


Пусть $(N,\ov{U\setminus E})$ --- орграф, полученный из графа 


$(N,U\setminus E)$, обратного 


к интервальному, с помощью некоторой транзитивной ориентации. 


Обозначим через $t(E)$ расписание, определенное соотношениями 


$$ 


t_i(E)=\max_{e\in L_i(\ov{U\setminus E})} 


v_0(e),\ \ i\in N, 


$$ 


где $L_i(\ov{U\setminus E})$ --- множество путей орграфа 


$(N,\ov{U\setminus E})$, для которых вершина 


$i$ является конечной. Пути заданы списками вершин (для вершин с нулевой 


полустепенью захода $L_i(\ov{U\setminus E})=\{i\}$). 


 


\begin{thm} 


Пусть $\bold P^0_\Omega=\{E^1,\dotsc,E^p\}$ --- полное множество 


альтернатив 


задачи 


\begin{equation} 


F(E)=(F_0(E),\dotsc,F_m(E))\to\min,\quad 


E\in\Omega_4=\Omega_0\cap\Omega_2. 


\end{equation} 


Тогда $\bold P^0_T=\{t(E^1),\dotsc,t(E^p)\}$. 


\end{thm} 


 


\begin{center} 


{\bf Алгоритм 2} ($F,g$-оракул задачи (4))


\end{center} 


 


1. Дано $E\in2^U$. Если $G(E)$ --- не интервальный граф, то 


$E\notin\Omega_4$; перейти на п.\,4. 


 


2. Проверить условия $F_j(E)\le\beta_j$, $j\in M$. Если хоть одно из 


них не 


выполняется, то $E\notin\Omega_4$; перейти на п.\,4. 


 


3. Получили $E\in\Omega_4$. Вычислить $F_0(E)$. 


 


4. Конец. 


%\smallskip 


 


 


Интервальность графа проверяется эффективно [8], поэтому 


релаксированная (относительно (2)) задача (4) имеет полиномиальную 


оракульную сложность. Нахождение эффективных расписаний, соответствующих 


эффективным траекториям, сводится к построению максимальных путей 


транзитивного орграфа. 


 


Можно доказать, что $\bold P_T=\mathop{\cup}\limits_{E\in\bold 


P'_\Omega}X^*(E)$, 


где $\bold P'_\Omega$ --- множество Парето 


задачи (4), $X^*(E)=\{t\in X(E):f_0(t)=F_0(E)\}\ne\emptyset$. Но, в 


отличие от (3), система 


множеств $X^*(E)$, $E\in\bold P'_\Omega$, является покрытием множества 


$\bold P_T$, т.\,к. подмножества, 


соответствующие разным траекториям, могут пересекаться. 


 


 


\begin{thebibliography}{9} 


 


\bibitem{1} 


{\em Теория расписаний и вычислительные машины} / Под. ред. Э.Г.\,Коффана. 


-- 


М.: Наука, 1984. -- 334 с. 


\bibitem{2} 


Гимади Э.Х., Залюбовский В.В., Шарыгин П.И. {\em Задача упаковки в 


полосу$:$ 


асимптотически точный подход}. // 


Изв. вузов. Математика. -- 1997. -- \No\,\,12. 


-- C.\,34--44. 


\bibitem{3} 


Зинченко А.Б. {\em Полиномиальная разрешимость специальных задач 


дизъюнктивного программирования} // Журн. вычисл. матем. и матем. физ. -- 


1999. -- Т.\,39. -- \No\,2. -- С.\,341--345. 


\bibitem{4} 


Козырев В.П. {\em Описание и порождение всех минимальных раскрасок 


интервального графа и решение смежных задач} // Журн. вычисл. матем. и 


матем. физ. -- 1996. -- Т.\,36. -- \No\,5. -- С.\,146--152. 


\bibitem{5} 


Кравцов М.К. {\em Неразрешимость задач векторной дискретной оптимизации 


в классе алгоритмов линейной свертки критериев} // Дискрет. матем. -- 1996.


-- Т.\,8. -- Вып.\,2. -- С.\,89--96. 


\bibitem{6} 


Подиновский В.В., Ногин В.Д. {\em Парето-оптимальные решения 


многокритериальных задач}. -- М.: Наука, 1982. -- 254 с. 


\bibitem{7} 


Peler Itsik, Shamir Ron. {\em Realizing interval graphs with size and


distance constraints} // SIAM J. Discrete Math. -- 1997. -- V.\,10. 


-- \No\,4. -- P.\,662--687. 


\bibitem{8} 


Korte N., Mohring R.H. {\em An incremental linear-time algorithm for 


recognizing interval graphs} // 


SIAM J. Comput. -- 1989. -- V.\,18. -- \No\,1. -- P.\,68--81. 


 


\end{thebibliography} 


 


\vskip 0.3cm 


 


\postup{\it Ростовский государственный }{\qquad \it Поступили} 


\postup{\it университет}{\it первый вариант $22.06.1999$} 


\postup{}{\it окончательный вариант $11.08.2000$} 


 


\label{end} 


\end{document} 





